% ===== PRÉAMBULE CANONIQUE — à coller VERBATIM en tête de CHAQUE .tex =====
\documentclass[11pt,a4paper]{article}
\usepackage[utf8]{inputenc}\usepackage[T1]{fontenc}\usepackage{lmodern}
\usepackage[french]{babel}
\usepackage{amsmath,amssymb,mathtools}\usepackage{icomma}
\usepackage{siunitx}\sisetup{locale=FR}  % \qty{}{}, \si{}, \ang{}
\usepackage{xcolor}\usepackage{colortbl}
\usepackage{tikz}\usepackage{pgfplots}\pgfplotsset{compat=1.18}
\usepackage[siunitx,europeanresistors]{circuitikz} % circuits (élec.)
\usepackage[most]{tcolorbox}\usepackage{enumitem}\usepackage{array,booktabs}
\usepackage[a4paper,margin=2.2cm]{geometry}
\usetikzlibrary{arrows.meta,calc,angles,quotes,positioning,patterns,%
  backgrounds,shapes.geometric,decorations.pathmorphing,decorations.markings,intersections}
\definecolor{primary}{RGB}{222,49,99}\definecolor{primarydark}{RGB}{150,22,55}
\definecolor{defcol}{RGB}{0,114,178}\definecolor{methcol}{RGB}{52,92,148}
\definecolor{excol}{RGB}{0,135,90}\definecolor{attncol}{RGB}{200,40,55}
\tcbset{boxbase/.style={enhanced,breakable,boxrule=0.9pt,arc=2.5pt,
  left=3mm,right=3mm,top=2.3mm,bottom=2.3mm,fonttitle=\bfseries,coltitle=white}}
% --- boîtes TUTORIEL (noms EXACTS, ne pas renommer) ---
\newtcolorbox{cle}[1][]{boxbase,colback=defcol!6,colframe=defcol,
  title={\textsc{À retenir}\ifx\relax#1\relax\relax\else\textmd{ : \itshape #1}\fi}}
\newtcolorbox{methode}[1][]{boxbase,colback=methcol!7,colframe=methcol,sharp corners,
  title={\textsc{Méthode}\ifx\relax#1\relax\relax\else\textmd{ --- \itshape #1}\fi}}
\newtcolorbox{exemplebox}[1][]{boxbase,colback=excol!5,colframe=excol,
  title={\textsc{Exemple}\ifx\relax#1\relax\relax\else\textmd{ : \itshape #1}\fi}}
\newtcolorbox{piege}[1][]{boxbase,colback=attncol!6,colframe=attncol,
  title={\textsc{Piège}\ifx\relax#1\relax\relax\else\textmd{ : \itshape #1}\fi}}
\newtcolorbox{remarque}[1][]{boxbase,colback=black!4,colframe=black!45,
  title={\textsc{Remarque}\ifx\relax#1\relax\relax\else\textmd{ : \itshape #1}\fi}}
% --- boîtes EXERCICE / CORRIGÉ (noms EXACTS, auto-comptées) ---
\newtcolorbox[auto counter]{exo}[1][]{boxbase,colback=primary!4,colframe=primary,
  title={\textsc{Exercice}~\thetcbcounter\ifx\relax#1\relax\relax\else\textmd{ --- \itshape #1}\fi}}
\newtcolorbox[auto counter]{corrige}[1][]{boxbase,colback=excol!4,colframe=excol,
  title={\textsc{Corrigé}~\thetcbcounter\ifx\relax#1\relax\relax\else\textmd{ --- \itshape #1}\fi}}
\newcommand{\vv}[1]{\vec{#1}}\newcommand{\ds}{\displaystyle}
\newcommand{\R}{\mathbb{R}}\newcommand{\N}{\mathbb{N}}\newcommand{\Z}{\mathbb{Z}}
% =========================================================================
\begin{document}
\usetikzlibrary{babel}
\begin{corrige}[filament : résistance à chaud et à froid]
\begin{enumerate}[label=\alph*)]
  \item Allumée : $R_{\text{on}}=\dfrac{U^{2}}{P}=\dfrac{200^{2}}{100}=\dfrac{40000}{100}=\qty{400}{\ohm}$.
  \item Éteinte : $R_{\text{off}}=\dfrac{R_{\text{on}}}{10}=\qty{40}{\ohm}$.
  \item Allumé, le filament est très chaud ; or la résistance d'un métal croît avec la température. La
        lampe est un dipôle \emph{non ohmique} : $R$ dépend de la température.
\end{enumerate}
\end{corrige}

\begin{corrige}[série ou parallèle : quel montage chauffe le plus ?]
\begin{enumerate}[label=\alph*)]
  \item \textbf{Série} : $R_{\text{eq}}=2R$, donc $W_1=\dfrac{U^{2}}{R_{\text{eq}}}\,t=\dfrac{U^{2}t}{2R}$.\;
        \textbf{Parallèle} : $R_{\text{eq}}=\dfrac{R}{2}$, donc $W_2=\dfrac{U^{2}}{R/2}\,t=\dfrac{2U^{2}t}{R}$.
  \item $\dfrac{W_2}{W_1}=\dfrac{2U^{2}t/R}{U^{2}t/(2R)}=4$. Le montage \textbf{parallèle} dissipe
        \textbf{4 fois plus} de chaleur (sa résistance équivalente est plus petite, donc l'intensité
        totale plus grande, et $W\propto I^2$).
\end{enumerate}
\end{corrige}

\begin{corrige}[chaleur dégagée par deux résistances en série]
\begin{enumerate}[label=\alph*)]
  \item $R_{\text{eq}}=R_1+R_2=\qty{60}{\ohm}$ ; \; $I=\dfrac{U}{R_{\text{eq}}}=\dfrac{20}{60}=\dfrac{1}{3}\,\si{\ampere}$.
  \item $t=9\times3600=\qty{32400}{\second}$, donc
        \[ W_{\text{th}}=R_{\text{eq}}I^{2}t=60\times\Bigl(\tfrac{1}{3}\Bigr)^{2}\times32400
           =60\times\tfrac{1}{9}\times32400=\qty{216000}{\joule}=\qty{216}{\kilo\joule}. \]
\end{enumerate}
\end{corrige}

\begin{corrige}[pourquoi transporter à haute tension ?]
\begin{enumerate}[label=\alph*)]
  \item $I_1=\dfrac{P}{U_1}=\dfrac{1000}{100}=\qty{10}{\ampere}$ ; pertes
        $P_{\text{p},1}=R_{\text{ligne}}I_1^{2}=1\times10^{2}=\qty{100}{\watt}$.
  \item $I_2=\dfrac{P}{U_2}=\dfrac{1000}{200}=\qty{5}{\ampere}$ ; pertes
        $P_{\text{p},2}=1\times5^{2}=\qty{25}{\watt}$. Les pertes ont été \textbf{divisées par 4}
        (tension $\times2$ $\Rightarrow$ intensité $\div2$ $\Rightarrow$ pertes $\div4$).
  \item \textbf{Faux} : les pertes valent $R_{\text{ligne}}I^{2}$ ; il faut au contraire \emph{diminuer}
        l'intensité (donc \emph{augmenter} la tension) pour les réduire.
\end{enumerate}
\end{corrige}

\end{document}
